% ── §7  Related Work ──────────────────────────────────────────────────────────
\section{Related Work}
\label{sec:related}

\paragraph{Tail Latency Mitigation.}
Dean and Barroso~\cite{dean2013tail} survey hedged requests and
tied requests as application-level mitigations. These are complementary
to \xtal, which operates at the infrastructure layer.
SEDA~\cite{welsh2001seda} uses staged event-driven queues with dynamic
thread-pool tuning but does not cross the streaming-to-storage boundary,
and predates modern ML-based control.

\paragraph{Streaming Engine Autoscaling.}
Naiad~\cite{murray2013naiad} and MillWheel~\cite{akidau2013millwheel}
address fault tolerance and consistency but do not expose closed-loop
latency control. Flink's Adaptive Scheduler~\cite{carbone2015flink}
adjusts task slot assignments reactively but does not tune checkpointing
or Kafka-side parameters, and does not produce explanations.
DS2~\cite{kalavri2018} and similar data-driven autoscalers optimise
throughput, not tail latency, and require offline profiling.

\paragraph{ML-Based System Tuning.}
CDBTune and OtterTune apply Bayesian optimisation and RL to database
knob tuning; their offline training phase is prohibitive for streaming
workloads with rapidly changing load. \xtal uses online contextual
bandits that adapt within minutes without offline profiling.
Decima~\cite{mao2019} applies RL to cluster job scheduling; it targets
job completion time, not per-event tail latency.

\paragraph{Explainable AI for Systems.}
SHAP has been applied to network traffic classification and AIOps alert
ranking, but not to control-plane decisions in streaming systems.
\xtal is the first to produce per-decision SHAP attributions in a
streaming-engine control loop.

\paragraph{Distributed Tracing.}
Dapper~\cite{sigelman2010dapper} and Jaeger make tail latency
\emph{diagnosable} post-hoc; \xtal makes it \emph{controllable}
in real time.
